% ==========================================================================
%  Chapter 100 -- Deflated Continuation (Extension Plan, §P1)
% ==========================================================================

\chapter{Deflated Continuation for Spurious-Root Avoidance}
\label{ch:deflated-continuation}

\paragraph{Normative status.}
Phase~P1 of the branch-selection extension
(Cap.~\ref{ch:branch-selection-reversals}). The deflation policy and its
solver integration are implemented in
\href{run:src/analysis/RegularizedNewtonContinuation.hh}{%
\texttt{RegularizedNewtonContinuation.hh}} and verified by ctest
\texttt{regularized\_newton\_deflation}
(\path{tests/test_regularized_newton_deflation.cpp}). Status:
\texttt{mechanism\_verified\_in\_isolation}; the default
(\texttt{NoDeflation}) path is compiled identically to the baseline, so
Gate~G1's non-regression clause holds by construction. The remaining
G1 clause --- bit-identical CSV against \texttt{g0\_baseline} with the
gate off, plus a converged real reversal --- is deferred to the driver
integration (Phase~I).

\section{Deflation operators}

Let \(\mathbf{R}:\mathbb{R}^n\to\mathbb{R}^n\) be the residual and
\(\mathbf{u}_1,\dots,\mathbf{u}_m\) a set of known (here: \emph{spurious})
roots. Following \cite{FarrellBirkissonFunke2015}, define the scalar
deflation weight
\[
  \eta(\mathbf{u})=\prod_{i=1}^{m}
      \Bigl(\lVert\mathbf{u}-\mathbf{u}_i\rVert^{-p}+\sigma\Bigr),
  \qquad p>0,\ \sigma\ge 0,
\]
and the deflated residual \(\mathbf{G}(\mathbf{u})
=\eta(\mathbf{u})\,\mathbf{R}(\mathbf{u})\). Two properties make
\(\mathbf{G}\) the right object:
%
\begin{enumerate}
  \item \emph{Root preservation.} Away from the deflated points,
        \(\eta>0\) is finite, so \(\mathbf{G}(\mathbf{u})=\mathbf{0}
        \iff \mathbf{R}(\mathbf{u})=\mathbf{0}\): every other equilibrium
        of the problem is preserved.
  \item \emph{Root repulsion.} As
        \(\mathbf{u}\to\mathbf{u}_i\), \(\eta\to\infty\) (for
        \(p>0\)); a Newton iteration on \(\mathbf{G}\) cannot re-converge
        to \(\mathbf{u}_i\). The shift \(\sigma\) keeps \(\eta\) bounded
        below away from the roots so the far-field problem is unchanged
        up to a positive scale.
\end{enumerate}
%
The exponent \(p\) sets how sharply the iteration is pushed off the
deflated root (\(p=2\) here); \(\sigma=1\) is the standard shift.

\section{Sherman--Morrison reduction of the deflated step}

Newton on \(\mathbf{G}\) requires \(\nabla\mathbf{G}\). Because
\(\mathbf{G}=\eta\mathbf{R}\),
\[
  \nabla\mathbf{G}
   = \eta\,\nabla\mathbf{R}
   + \mathbf{R}\,(\nabla\eta)^{\mathsf T},
\]
a \emph{rank-one} update of the scaled tangent \(\eta\nabla\mathbf{R}\).
The deflated Newton step solves
\(\nabla\mathbf{G}\,\delta\mathbf{u}_G=-\mathbf{G}\). Applying the
Sherman--Morrison identity to the rank-one update, and writing the
\emph{ordinary} (undeflated) Newton direction
\(\delta\mathbf{u}_N=-(\nabla\mathbf{R})^{-1}\mathbf{R}\), the deflated
direction collapses to a scalar rescaling
\cite{FarrellBirkissonFunke2015}:
\[
  \boxed{\;\delta\mathbf{u}_G=\tau\,\delta\mathbf{u}_N,
  \qquad
  \tau=\frac{1}{\,1-\dfrac{\nabla\eta\cdot\delta\mathbf{u}_N}{\eta}\,}.\;}
\]
The whole cost of deflation is therefore \emph{one extra scalar} per
Newton iteration, assembled from cheap vector operations. With
\(\nabla\eta/\eta=\sum_i\nabla m_i/m_i\) and
\(m_i=\lVert\mathbf{u}-\mathbf{u}_i\rVert^{-p}+\sigma\),
\[
  \frac{\nabla\eta}{\eta}\cdot\delta\mathbf{u}_N
   =\sum_{i=1}^{m}
      \frac{-p\,\lVert\mathbf{u}-\mathbf{u}_i\rVert^{-p-2}\,
            \bigl(\mathbf{u}-\mathbf{u}_i\bigr)\cdot\delta\mathbf{u}_N}
           {\lVert\mathbf{u}-\mathbf{u}_i\rVert^{-p}+\sigma},
\]
i.e.\ per registered root one \texttt{VecWAXPY} (form
\(\mathbf{u}-\mathbf{u}_i\)), one \texttt{VecNorm} and one
\texttt{VecDot}. The scalar \(\tau\) is clamped to
\(\pm\tau_{\max}\) (sign preserved) so a near-singular denominator cannot
produce an unbounded step.

\section{Interaction with Levenberg--Marquardt: the acceptance test}

The subtlety of combining deflation with LM is the \emph{acceptance
test}. The baseline solver accepts a trial iterate only if the
\emph{undeflated} residual decreases, \(\lVert\mathbf{R}(\mathbf{u}
+\delta\mathbf{u})\rVert<\lVert\mathbf{R}(\mathbf{u})\rVert\). But a step
that \emph{escapes} the spurious basin must be allowed to raise
\(\lVert\mathbf{R}\rVert\) locally while it climbs out. The fix is to
monitor the \emph{deflated} residual across the step:
\[
  \text{accept}\iff
  \lVert\mathbf{R}(\mathbf{u}+\delta\mathbf{u})\rVert\,
    \eta(\mathbf{u}+\delta\mathbf{u})
  \;<\;
  \lVert\mathbf{R}(\mathbf{u})\rVert\,\eta(\mathbf{u}).
\]
Near a registered root \(\eta\) is large, so the deflated measure
penalises staying put and rewards the move away. Convergence and the
\texttt{accept\_floor} continue to be judged on the \emph{pure}
\(\lVert\mathbf{R}\rVert\), so the deflation changes only \emph{which}
step is taken, never the definition of a solution. When no roots are
registered, \(\eta\equiv1\) and the test reduces to the baseline
comparison.

\section{Implementation}

The mechanism is a compile-time policy, so the default solver is
untouched:
%
\begin{itemize}
  \item \texttt{concept DeflationPolicy} with two models: \texttt{NoDeflation}
        (\texttt{enabled=false}) and
        \texttt{ShermanMorrisonDeflation\{power,shift,tau\_max\}}.
  \item A third class template parameter
        \texttt{Defl = NoDeflation}. \emph{All} deflation code lives
        behind \texttt{if constexpr (Defl::enabled)}; when disabled the
        emitted object code is identical to the pre-P1 solver.
  \item In \texttt{advance\_to}: after the \texttt{KSPSolve} the Newton
        direction is scaled by \(\tau\); the acceptance test uses the
        deflated measure above.
  \item A branch-selection API for the driver's detect--restore--retry
        loop: \texttt{register\_spurious\_root},
        \texttt{clear\_spurious\_roots}, \texttt{set\_solution} (which
        also resets the secant history), \texttt{set\_config},
        \texttt{set\_regularization}, and
        \texttt{capture\_state}/\texttt{restore\_state} for the cheap
        checkpoint that brackets every retry (a spurious step is detected
        only after the crack ratchet has advanced, hence the rollback).
\end{itemize}

\section{Verification (ctest \texttt{regularized\_newton\_deflation})}

The toy backend is the 1-DOF cubic \(R(u)=(u-1)(u-2)(u-4)\), whose stable
roots (\(R'>0\)) are \(u=1\) and \(u=4\) and whose unstable root is
\(u=2\). Three checks:
%
\begin{description}
  \item[A --- baseline attraction.] From \(u_0=3.8\) the undeflated LM
        continuation lands on the high stable root \(u=4\)
        (\(|R|=6.3\times10^{-10}\), converged).
  \item[B --- deflated escape.] With \(u=4\) registered as spurious, the
        deflated solver from the same \(u_0=3.8\) is repelled and
        converges to a \emph{different} root, \(u=2\)
        (\(|R|=8.2\times10^{-9}\), converged). This is the branch-selection
        mechanism working end to end.
  \item[C --- rollback.] \texttt{capture\_state}/\texttt{restore\_state}
        round-trips the solution to machine precision after an
        intervening perturbed solve.
\end{description}
%
The test passes with zero failures. It deliberately does \emph{not}
assert the pre-P1 anecdote ``\(u_0=3.5\to\) high root'': at \(u_0=3.5\)
the ordinary Newton step overshoots root~\(4\) to \(u\approx4.5\), where
the residual is larger, and the monotone-residual LM acceptance
\emph{rejects} the overshoot and stalls. That is correct solver
behaviour --- the plan's anecdote assumed pure Newton (which overshoots
then returns), whereas this solver's guarded acceptance does not --- and
it is why the isolated test seeds from \(u_0=3.8\), where the baseline
converges cleanly and the deflation contrast is unambiguous.

\section{Honest limits}

The scalar \(\tau\) is exact for a pure Newton step; combined with a
large LM \(\mu\) the correction is approximate, which is why the deflated
\emph{acceptance} test (not just the direction scaling) is essential and
why \(\tau_{\max}\) caps the rescaling. If a reversal lives pinned at
\(\mu_{\max}\), deflation may not suffice on its own --- that is the
regime the energy remedy (Cap.~\ref{ch:energy-tao-continuation}) and the
meta-tuner (Cap.~\ref{ch:cultural-algorithm-module}) are meant to cover.
The promotion decision between the three remedies is a Phase-I evidence
matrix, not a foregone conclusion.
